\begin{table}[!htbp]\centering
\caption{Income's incremental contribution among high-math students}
\label{tab:high-math-main-cubic-income-contribution}
\small
\begin{tabular}{lrrrr}\hline\hline
Outcome & $R^2$ without income & Full $R^2$ & Income $\Delta R^2$ & Income share (\%) \\ \hline
Premium institution & 0.1108 & 0.1667 & 0.0559 & 33.54 \\
Premium field & 0.0573 & 0.0756 & 0.0182 & 24.11 \\
Either (or both) & 0.1287 & 0.1308 & 0.0021 & 1.62 \\
Both & 0.0927 & 0.0992 & 0.0065 & 6.53 \\
\hline\hline\end{tabular}
\begin{minipage}{0.98\textwidth}\footnotesize
\textit{Notes:} Income share is the increase in ordinary unadjusted $R^2$ from adding nine income-decile indicators, divided by full-model $R^2$. Both models retain all nine total-degree-three math/verbal terms and admission-year indicators, using the same 93,423 students. Selection requires original $z_{math} \geq 1$. This is the unique incremental contribution after score/year controls; shared explanatory variation is not allocated to income. The 840 admission-2026 records remain. All fit measures are in-sample and descriptive.
\end{minipage}
\end{table}
